\section{Ausreißer machen das Ergebnis unschaaf - mäh}
\textcolor{red}{TODO: weitere T-Tests von Aufgaben 2+3}\\
?? hier? oder war das bei der Nutzerbefragung?

% ****************************************************************
%   Effektivität
% ****************************************************************
Aufgaben angemessenheit
wenn bei 2 die beiden Ausreißer weggenommen werden, die andersrum bewertet haben, dann wird der Unterschied auch signifikant (s. \autoref{fig:hist_prümper_all})\\
Boxplot: \autoref{fig:pruemper_Aufgabenangemessenheit}

Steuerbarkeit
wenn bei 3 alt die beiden Ausreißer weggenommen werden, die andersrum bewertet haben, dann wird der Unterschied auch signifikant (s. \autoref{fig:hist_prümper_all})\\
Boxplot: \autoref{fig:pruemper_Steuerbarkeit}

-	Die Hypothesen die nicht bestätigt werden können
o	gut für mich
 können nicht bestätigte Hypothesen durch ungünstige Fälle doch bestätigt werden?
 wird das Ergebnis noch signifikanter?
o	schlecht für mich
 muss eine Gegendarstellung betrachtet werden?
 wo führen Extrema zur Signifikanz, die sonst nicht da ist?
-	Beispiele
o	neuer T-Test weil Ausreißer, ist dann doch signifikant